\documentclass[11pt]{article}
\usepackage[margin=0.84in]{geometry}\usepackage{amsmath,amssymb,amsthm,microtype,hyperref}
\hypersetup{colorlinks=true,linkcolor=blue,urlcolor=blue}
\newtheorem{theorem}{Theorem}\newtheorem{lemma}[theorem]{Lemma}
\title{Recovering the modulus and quadratic twist\\from one finite band of signed zeros}
\author{Research note on the Dirichlet GRH development}\date{}
\begin{document}\maketitle
\begin{abstract}
The signed zero multiplicities in the Lean proof forget the names of the
characters and allow cancellations at common zeros. An earlier ProveIt note
recovered the parameters from limiting high-band data and asked for a
certified stopping rule. We give one for the lattice parameters, with
explicit constants and an error model. Its conductor bound is inferred from
the same band used for reconstruction. The estimates prioritize certification over numerical efficiency.
\end{abstract}
\section{Data and conclusion}
Suppose $18480\mid q$ and $d$ is a positive prime with
$d\equiv q/2-1\pmod q$. Then $(d,q)=1$, $8\mid q$, $d\equiv7\pmod8$,
and $q\le 2(d+1)$. Let $h=\varphi(q)$, $\eta(n)=(n/d)$, and let
$\chi^*$ denote the primitive character inducing $\chi$. For
$\rho=1/2+i\gamma$ define
\[
 m(\gamma)=h^{-1}\sum_{\chi\bmod q}
 \bigl(\operatorname{ord}_\rho L(s,\chi^*)-
       \operatorname{ord}_\rho L(s,\chi^*\eta)\bigr).
\]
These are the signed orders in the residue calculation for the Lean
development's smoothed prime remainder \cite{lean}. Dirichlet GRH
accounts for all the nontrivial zeros. Both signs of the ordinate are used.
\begin{theorem}
From a sufficiently high band $T\le |\gamma|\le 2T$ of the aggregate
signed data one can recover $(q,d)$ with a computable certificate.
No upper bound for $q$ or $d$ is part of the input. On an insufficient
band the procedure may return \emph{undetermined}. If the band measurements
defined below have certified absolute errors $e_u(T)$ tending to zero for
each fixed observation site $u$, the procedure succeeds on every
sufficiently high band.
\end{theorem}
A single accepted band certifies both integers. This assumes a genuine packet
in the stated parameter class. Authenticating a putative spectrum is a separate problem.
\section{Identities and a conductor bound}
Write
\[
 F(t)=\sum_\gamma\frac{m(\gamma)e^{i\gamma t}}{1/4+\gamma^2},
 \qquad \mu=(D^2-1/4)F.
\]
The series for $F$ is absolutely convergent. As a distribution, $\mu$
has an atom of mass $J(0)=\log d$, smooth density
$-2\cosh(t/2)/h$, and remaining atoms among $\pm\log(p^k)$.
At a prime $p\nmid qd$ its mass is
\begin{equation}
 J(\log p)=\frac{\log p}{\sqrt p}(1-\eta(p))\mathbf1_{p\equiv1\bmod q}.
 \label{eq:mass}
\end{equation}
These identities were proved in the archived \texttt{cusp.tex} \cite{incoming}.
Apply the primitive Dirichlet explicit formula to both
families \cite{chirre}. The gamma terms cancel because the two families
have identical parity counts. The conductor is multiplied by $d$ in every
twist. Only the untwisted principal character produces poles, contributing
the density above after division by $h$. At unramified prime powers average
the character values by orthogonality and multiply by $1-\eta(p^k)$.
At ramified primes the coefficients change but the support
still consists of logarithms of prime powers. The sign follows from
$\mu=-\sum m(\gamma)e^{i\gamma t}$.

\begin{lemma}
For this entire parameter class, $\|F\|_\infty\le 9\log d+60$.
\end{lemma}
\begin{proof}
For a primitive character of conductor $f$, let $\Xi$ be its completed
function, with an additional factor $s(s-1)$ in the zeta case. Its
Hadamard logarithmic derivative and functional equation give
\[
 \sum_\rho \Re\frac1{2-\rho}=\Re\frac{\Xi'}{\Xi}(2)
 \le \tfrac12\log f+4.
\]
Indeed the constant in the real part of the Hadamard logarithmic
derivative vanishes by evaluating at a regular point on the critical
line, where both real parts vanish by GRH and reflection. At $2$ the
gamma contribution is at most $1/2$, since the digamma function is
increasing and $\psi(2)<1$. The logarithmic Euler derivative has absolute
value at most $\sum_{n\ge2}(\log n)/n^2<2$ by the integral test, and the
zeta correction is $1+1/2$. Discard $-\tfrac12\log\pi$. On GRH,
$(1/4+\gamma^2)^{-1}\le 6\Re(2-\rho)^{-1}$. Adding the two families
therefore gives $\|F\|_\infty\le6\log q+3\log d+48$.
Now $q\le 2(d+1)\le4d$.
\end{proof}

\section{A computable bound from a high band}
Use $\widehat g(s)=\int g(t)e^{-ist}\,dt$. Set
\[
 g(t)=\cos(3t/2)\bigg(\frac{\sin(t/16)}{t/16}\bigg)^8,\qquad g(0)=1.
\]
Its transform is supported on $[-2,-1]\cup[1,2]$. Each sinc factor
has frequency support $[-1/16,1/16]$. The transform is an explicit
piecewise polynomial, obtained by eight convolutions of interval
indicators followed by two translations. Define
\[
 A_T(u)=-T^{-1}\sum_{T\le|\gamma|\le2T}m(\gamma)
                \widehat g(\gamma/T)e^{i\gamma u}.
\]
It equals $\langle\mu,g(T(\cdot-u))\rangle$, with the pairing defined as
$\int F(t)(D^2-1/4)[g(T(t-u))]\,dt$.
Termwise integration follows from absolute uniform convergence of $F$.

Choose $0<r\le1$ such that $(u-2r,u+2r)$ contains no potential atom except
possibly $u$. At zero use $r=\log(2)/3$. At $u=\log p$ take one third
of $\min\{\log(p/(p-1)),\log((p+1)/p)\}$. For $Tr\ge1$, put
\begin{align*}
 B_u(T)&=\frac{512}{7T}\cosh\big((|u|+2r)/2\big),\\
 K_u(T)&=\frac{2\cdot16^8}{7r^7 T^8}
        \left(108T^2+40T/r+6/r^2+1/4\right).
\end{align*}
Then
\begin{equation}
 |A_T(u)-J(u)|\le B_u(T)+(9\log d+60)K_u(T).
 \label{eq:error}
\end{equation}
To verify the constants, multiply the test by a cutoff equal to one on
$|t-u|\le r$, zero on $|t-u|\ge2r$, and between these use
$1-10z^3+15z^4-6z^5$ at $z=|t-u|/r-1$.
Its first two derivative bounds are $2/r,6/r^2$. On the local part only
the selected mass and the density contribute. Since
$\int|g|\le2(16+16/7)=256/7$, the latter costs at most $B_u(T)$.
For the complementary part integrate by parts onto the test. Writing
$g(t)=16^8t^{-8}\cos(3t/2)\sin^8(t/16)$ shows that for $|t|\ge1$
its derivatives of orders $0,1,2$ are bounded by
$16^8|t|^{-8}$ times $1,10,108$, respectively. The trigonometric
factor has frequency at most $2$ and total absolute exponential coefficients
at most $1$. The product rule and
$\int_{|t-u|\ge r}|t-u|^{-8}dt=2/(7r^7)$ give $K_u$.
The cutoff is $C^2$, which suffices by approximation. The integrable
differentiated tests justify passage to the full real line.

\section{Certification and noise}
Let $a$ be the real part of the observed $A_T(0)$ and write $K=K_0(T)$,
$D=B_0(T)+60K+e_0(T)$. When $9K<1$, \eqref{eq:error} confines $\log d$ to
\[
 \left[\frac{a-D}{1+9K},\frac{a+D}{1-9K}\right].
\]
Use outward-rounded arithmetic and test the primes in the exponential of
this interval. If there is a unique eligible prime, it is $d$. Otherwise
decline this band. For fixed true $d$ the interval shrinks to $\log d$.
List the divisors $q'$ of $2(d+1)$ with $18480\mid q'$ and
$d\equiv q'/2-1\pmod {q'}$.

For every pair of distinct candidate moduli, find a prime $p$ coprime to
all candidates and $d$, with $\eta(p)=-1$, for which their
indicators in \eqref{eq:mass} differ. Enumerating primes terminates.
Indeed choose a prime power $\ell^a$ occurring to a higher exponent in one
modulus than in the other. Modulo the full least common multiple of the
candidates choose $1$ at all other components and a unit congruent to $1$
modulo the lower $\ell$-power but distinct from $1$ modulo $\ell^a$.
If the lower power is $1$, then $\ell$ is odd, since $8$ divides each
candidate. Set a nonresidue at $d$ independently. Dirichlet's theorem
supplies infinitely many primes in this class.
For each witness intersect the two possible masses, $0$ and
$2\log p/\sqrt p$, with the certified interval around the observed
$\Re A_T(\log p)$ from \eqref{eq:error} plus measurement error. Retain
compatible moduli. Each witness eventually separates its pair because $p$
is fixed and $B_u,K_u,e_u$ tend to zero. A surviving singleton is certified.

The errors are errors in individual band readings. For example, let
$E_T$ bound the total absolute missing, inserted, or altered weight and
let $\delta_T$ bound the displacement of matched ordinates. Use records
in a padded band so every entry that can cross an endpoint is included.
If $M_T$ bounds the absolute weight of the matched entries, a valid error is
\[
 e_u(T)=\frac{\|\widehat g\|_\infty E_T}{T}
 +\frac{M_T\delta_T}{T}
 \left(|u|\|\widehat g\|_\infty+\|\widehat g'\|_\infty/T\right)
 +e_{\rm arithmetic}.
\]
Both norms can be bounded from the explicit spline. For each fixed packet
$M_T=O_{q,d}(T\log(2+T))$ by zero counting. Thus
$E_T=o(T)$ and $\delta_T=o(1/\log T)$ suffice asymptotically for a
fixed-width padding when $\delta_T\le1$, if the
arithmetic error tends to zero. For actual certificates $M_T$ can be
bounded by the observed absolute weight plus the certified weight budget.
The earlier single-band proof already establishes limiting recovery
and robustness. The affine bound in \eqref{eq:error} supplies the
test for accepting a band before either parameter is known.
\paragraph{A finite example.}
Take $d=46199$ and either $q=18480$ or $q=92400$. A genuine band at
$T=10^{12}$ suffices, even with absolute reading error $10^{-7}$ at each
of the two sites. Use $r=1/8$ at zero and $r=1/(4p)$ at $p=110881$.
The origin interval contains only $d$, which leaves exactly these two
candidate moduli. The witness $p$ is an inert prime, with $p-1=6\cdot18480$.
The two possible masses differ by more than $2/334$. The accompanying
\texttt{check\_decoder\_example.py} verifies the error margins by rational
arithmetic. This is an a priori guarantee for either packet; no zero
ordinate at height $10^{12}$ was computed.
\begin{thebibliography}{9}
\bibitem{lean} L. Sernau and collaborators,
\href{https://github.com/lsernau/grh-lean/tree/72fead102f663d5c7a407964a13dc0cf9d977207}{\emph{GRH for Dirichlet $L$-functions, in Lean 4}},
especially \texttt{grh/src/New/Reduction/Zeros}.
\bibitem{incoming} V. Reshetnikov, ProveIt incoming collection,
\href{https://github.com/VladimirReshetnikov/ProveIt/blob/d31039a48d963692e5a5176f47b6603fc9e9d9cf/docs/incoming/dirichlet-grh-research.zip}{\texttt{dirichlet-grh-research.zip}},
\texttt{critical-line/cusp.tex}, \texttt{band.tex}, and research question 3
in \texttt{critical-line.tex}.
\bibitem{chirre} A. Chirre, V. J. Pereira J\'unior and D. de Laat,
\href{https://doi.org/10.1090/mcom/3638}{Primes in arithmetic progressions
and semidefinite programming}, Math. Comp. 90 (2021), section 2.1.
\end{thebibliography}\end{document}
